
\section{\BookRecordIndex{T0221}黎曼--罗赫定理的证明}\label{sec:C-4}

定理~\ref{thm:C-1-10} 的第一部分已由定理~\ref{thm:C-2-3} 证明。
因此只须建立指标公式，并证明 (iii) 中线丛的单射性与满射性结论。
注意，证明指标公式时只须考虑光滑复线性算子：
由引理~\ref{lem:C-2-2}，每个 $W^{\ell-1,p}$ 类\BookRecordIndex{T0210}柯西--黎曼算子
与这种算子相差一个紧算子，故由定理~\ref{thm:A-1-5}，两者指标相同。
又由定理~\ref{thm:C-1-10}(i)，只须在 $k=1$、$q=2$ 时证明指标公式。
然而 (iii) 的断言更为精细，必须研究给定的特定算子。

\begin{theorem}\label{thm:C-4-1}
若 $\Sigma=\mathbb D$ 是 $\mathbb C$ 中的闭单位圆盘，
且算子 $D$ 复线性，则定理~\ref{thm:C-1-10} 成立。
\end{theorem}

\begin{proof}
由\BookRecordIndex{T0062}边界马斯洛夫指标的直和公理及\BookRecordIndex{T0129}弗雷德霍姆指标相应的显然性质，
可设 $E\to\mathbb D$ 是复线丛。由推论~\ref{cor:C-3-9}，还可设
\[
  E=\mathbb D\times\mathbb C,\qquad
  F_{e^{i\theta}}=\mathbb R e^{ik\theta/2}
\]
对每个 $\theta\in\mathbb R$ 和某个整数 $k$ 成立。证明分四步。
设
\[
  X_F:=W_F^{1,2}(\mathbb D,\mathbb C),\qquad
  Y:=L^2(\mathbb D,\Lambda^{0,1}T^*\mathbb D\otimes\mathbb C),
\]
并记 $D_F:X_F\to Y$ 为由此子丛及零联络给出的\BookRecordIndex{T0210}柯西--黎曼算子：
\[
  D_F\xi=\frac12(\partial_s\xi+i\partial_t\xi)(ds-i\,dt).
\]
前三步直接计算 $D_F$，最后一步将结果推广到其他算子。

\noindent\textbf{步骤 1：} $\operatorname{im}D_F$ 的正交补是所有形如
$\zeta\,d\bar z$ 的 $(0,1)$ 形式，其中
$\zeta:\mathbb D\to\mathbb C$ 光滑且满足
\begin{equation}\label{eq:C-4-1}
  \partial_s\zeta-i\partial_t\zeta=0,\qquad
  \zeta(e^{i\theta})\in i e^{i\theta+ik\theta/2}\mathbb R.
\end{equation}
关于正交补中的形式光滑且满足第一式，
可由式~\eqref{eq:C-1-3} 及注~\ref{rem:C-1-4} 中的对偶算子公式\BookRecordIndex{T0326}推出。
现取 $\xi\in X_F$ 及满足 $\partial_s\zeta-i\partial_t\zeta=0$
的 $\zeta:\mathbb D\to\mathbb C$。则
\begin{align*}
  \int_{\mathbb D}\langle\zeta\,d\bar z,D_F\xi\rangle\,ds\,dt
  &=\operatorname{Re}\int_{\mathbb D}
      \bar\zeta(\partial_s\xi+i\partial_t\xi)\,ds\,dt
    +\operatorname{Re}\int_{\mathbb D}
      \overline{(\partial_s\zeta-i\partial_t\zeta)}\,\xi\,ds\,dt\\
  &=\operatorname{Re}\int_{\mathbb D}
      \bigl(\partial_s(\bar\zeta\xi)+i\partial_t(\bar\zeta\xi)\bigr)\,ds\,dt\\
  &=\operatorname{Re}\int_0^{2\pi}
      e^{i\theta}\bar\zeta(e^{i\theta})\xi(e^{i\theta})\,d\theta.
\end{align*}
右端对所有 $\xi\in X_F$ 消失，当且仅当 $\zeta$ 满足式~\eqref{eq:C-4-1}。

\noindent\textbf{步骤 2：} 若 $k\le-1$，则 $D_F$ 单射；
若 $k\ge0$，则 $\dim\ker D_F=1+k$。
若 $\xi\in X_F$ 且 $\bar\partial\xi=0$，则 $\xi$ 有幂级数展开
\[
  \xi(z)=\sum_{n\ge0}a_nz^n,
\]
其收敛半径至少为 $1$。边界条件为
$\xi(e^{i\theta})e^{-ik\theta/2}\in\mathbb R$，等价于
\[
  \xi(e^{i\theta})e^{-ik\theta}
  =\overline{\xi(e^{i\theta})},\qquad\theta\in\mathbb R.
\]
Fourier 系数为
\[
  a_n=\lim_{r\to1}\frac1{2\pi}\int_0^{2\pi}
    \xi(re^{i\theta})e^{-in\theta}\,d\theta
  =\frac1{2\pi}\int_0^{2\pi}
    \xi(e^{i\theta})e^{-in\theta}\,d\theta.
\]
故
\[
  a_{k-n}=\frac1{2\pi}\int_0^{2\pi}
    \xi(e^{i\theta})e^{-ik\theta}e^{in\theta}\,d\theta
  =\frac1{2\pi}\int_0^{2\pi}
    \overline{\xi(e^{i\theta})}e^{in\theta}\,d\theta
  =\bar a_n.
\]
因此 $k<0$ 时 $\xi=0$；而 $k\ge0$ 时
\[
  \xi(z)=a_0+a_1z+\cdots+a_kz^k,\qquad a_{k-n}=\bar a_n.
\]
于是 $k\ge0$ 时 $\dim\ker D_F=1+k$。

\noindent\textbf{步骤 3：} 若 $k\ge-1$，则 $D_F$ 满射；
若 $k\le-2$，则 $\dim\operatorname{coker}D_F=-k-1$。
暂把上述带边界条件 $F(e^{i\theta})=\mathbb R e^{ik\theta/2}$
的算子记作 $D_F(k)$。若 $\zeta$ 满足式~\eqref{eq:C-4-1}，
则 $\bar\zeta\in\ker D_F(-k-2)$。因此由步骤 1，
$D_F(k)$ 满射恰当且仅当 $D_F(-k-2)$ 单射。
由步骤 2，当 $-k-2\le-1$，即 $k\ge-1$ 时满足这一条件，
证明了第一项；第二项留给读者。

\noindent\textbf{步骤 4：} 证明定理。圆盘上任意复线性\BookRecordIndex{T0210}柯西--黎曼算子
均有形式 $D_F+\alpha$，其中
$\alpha\in W^{\ell-1,p}(\mathbb D,\Lambda^{0,1}T^*\mathbb D)$。
由步骤 3 在 $k=0$ 的结论，存在
$f\in W^{\ell,p}(\mathbb D,\mathbb C)$，使
$f(e^{i\theta})\in\mathbb R$ 且 $\bar\partial f=\alpha$。
于是 $w:=e^f:\mathbb D\to\mathbb C^*$ 满足
\[
  wF=F,\qquad w^{-1}\bar\partial w=\alpha.
\]
故 $w^{-1}\circ D_F\circ w=D_F+\alpha$。从而
$D_F+\alpha$ 的单射性或满射性恰与 $D_F$ 相同。
这完成了定理~\ref{thm:C-4-1} 的证明。
\end{proof}

下面的\BookRecordIndex{T0221}黎曼--罗赫定理证明，将使用 C.3 节说明的裤形归纳。
先建立归纳步骤。设
\[
  \Sigma_{02}=\Sigma_{01}\cup\Sigma_{12},\qquad
  \Sigma_{01}\cap\Sigma_{12}=\partial\Sigma_{01}\cap\partial\Sigma_{12}
\]
是定义~\ref{def:C-3-1} 中的分解。记
\[
  \partial\Sigma_{ij}=(-\Gamma_i)\cup\Gamma_j,
\]
其中 $\Gamma_i\cap\Gamma_j=\varnothing$，且
$\Gamma_1=\Sigma_{01}\cap\Sigma_{12}$。
此符号表示：$\Gamma_1$ 取 $\Sigma_{01}$ 的边界定向；
$\Gamma_2$ 取 $\Sigma_{02}$ 或 $\Sigma_{12}$ 的边界定向；
$\Gamma_0$ 的定向与作为 $\Sigma_{01}$ 或 $\Sigma_{02}$ 边界时的定向相反。
设 $\Sigma_{02}$ 上的\BookRecordIndex{T0080}丛对 $(E_{02},F_{02})$ 如定义~\ref{def:C-3-4}
分解为 $(E_{01},F_{01})$ 与 $(E_{12},F_{12})$，其中
$F_{ij}=F_i\cup F_j$，$F_i\subset E_{02}|_{\Gamma_i}$ 是全实子丛。记
\[
  X_{ij}:=W_{F_{ij}}^{1,2}(\Sigma_{ij},E_{ij}),\qquad
  Y_{ij}:=L^2(\Sigma_{ij},\Lambda^{0,1}T^*\Sigma_{ij}\otimes_{\mathbb C}E_{ij}),
\]
并设 $D_{ij}:X_{ij}\to Y_{ij}$ 是限制 $E_{02}$ 上某个光滑\BookRecordIndex{T0210}柯西--黎曼算子所得的算子。

\begin{theorem}\label{thm:C-4-2}
在上述情形，\BookRecordIndex{T0129}弗雷德霍姆指标满足
\begin{equation}\label{eq:C-4-2}
  \operatorname{index}(D_{02})
  =\operatorname{index}(D_{01})+\operatorname{index}(D_{12}).
\end{equation}
\end{theorem}

\begin{proof}
证明的基本想法相当直接。设 $E_1\to\Gamma_1$ 是 $E_{02}$
在 $\Gamma_1$ 上的限制，$\bar E_1\to\Gamma_1$ 是把 $E_1$
的复结构反转所得的丛。则 $\bar E_1\oplus E_1$ 中的全实子丛
$F_1\oplus F_1$，在任意平凡化下与对角线有相同的马斯洛夫指标，
因而二者可经全实子丛同伦相连。该同伦中的每个全实子丛，
都为\BookRecordIndex{T0210}柯西--黎曼算子给出非局部椭圆边界条件；结论遂由
\BookRecordIndex{T0129}弗雷德霍姆指标的形变不变性\BookRecordIndex{T0326}推出。
为使此思路精确，将边界条件的同伦转化为固定算子之低阶扰动的同伦更方便。
为便于叙述，先设 $\Gamma_1$ 连通；一般情形只须对 $\Gamma_1$
各分支重复下面的构造。证明分五步；前两步将给定算子在 $\Gamma_1$ 附近标准化。

\noindent\textbf{步骤 1：} 取 $\Gamma_1$ 的闭管状邻域 $U\subset\Sigma_{02}$，
并选微分同胚
\[
  \phi:[-1,1]\times\mathbb R/\mathbb Z\longrightarrow U,
\]
使 $\phi(0\times\mathbb R/\mathbb Z)=\Gamma_1$。
由定理~\ref{thm:C-2-3}，不失一般性可设 $\phi$ 全纯。
于是 $i=\phi^*j$ 是坐标 $s+it$ 下
$[-1,1]\times\mathbb R/\mathbb Z$ 上的标准复结构，
其中 $-1\le s\le1$、$t\in\mathbb R/\mathbb Z$。

\noindent\textbf{步骤 2：} 选复平凡化
$U\times\mathbb C^n\to E|_U:(z,\zeta)\mapsto\Phi(z)\zeta$。
我们将证明：不改变指标即可使 $D\circ\Phi=\Phi\circ\bar\partial$。
\sourcefootnote{原书在定义 $A$ 之前便称可假定此等式成立；
一般的 $D$ 尚有下文定义的零阶项 $A$，该等式是以 $D-B$ 代替 $D$
之后才得到的。此处依原书随后给出的构造调整叙述次序。}
对常值映射 $\zeta_0\in\mathbb C^n$，以
$A\zeta_0:=\Phi^{-1}D(\Phi\zeta_0)$ 定义
$A\in\Omega^{0,1}(U,\operatorname{End}_{\mathbb R}(\mathbb C^n))$。
于是对每个光滑映射 $\zeta:U\to\mathbb C^n$，有
\[
  D(\Phi\zeta)=\Phi(\bar\partial\zeta+A\zeta).
\]
将 $\Phi A\Phi^{-1}\in\Omega^{0,1}(U,\operatorname{End}_{\mathbb R}(E))$
延拓为全局的自同态值 $(0,1)$ 形式
$B\in\Omega^{0,1}(\Sigma,\operatorname{End}_{\mathbb R}(E))$，
并以指标相同的\BookRecordIndex{T0210}柯西--黎曼算子 $D-B$ 代替 $D$。
这就证明了步骤 2。

\noindent\textbf{步骤 3：} 定义空间
\begin{align*}
  X&:=W^{1,2}(\Sigma_{01},E_{01})\oplus
      W^{1,2}(\Sigma_{12},E_{12}),\\
  Y&:=L^2(\Sigma_{01},\Lambda^{0,1}T^*\Sigma_{01}\otimes E_{01})
      \oplus L^2(\Sigma_{12},\Lambda^{0,1}T^*\Sigma_{12}\otimes E_{12}),
\end{align*}
以及子空间
\begin{align*}
  X_0:=\{(\xi_{01},\xi_{12})\in X\mid&
     \ \xi_{01}(\Gamma_0)\subset F_0,
       \ \xi_{01}(\Gamma_1)\subset F_1,\\
    &\ \xi_{12}(\Gamma_1)\subset F_1,
       \ \xi_{12}(\Gamma_2)\subset F_2\},\\
  X_1:=\{(\xi_{01},\xi_{12})\in X\mid&
     \ \xi_{01}(\Gamma_0)\subset F_0,
       \ \xi_{12}(\Gamma_2)\subset F_2,\\
    &\ \xi_{01}(z)=\xi_{12}(z),\ z\in\Gamma_1\}.
\end{align*}
（参见定义~\ref{def:C-1-5} 前的注。）
于是 $D$ 所确定的算子 $D_0:X_0\to Y$ 和 $D_1:X_1\to Y$
是\BookRecordIndex{T0128}弗雷德霍姆算子，且
\[
  \operatorname{index}(D_0)
    =\operatorname{index}(D_{01})+\operatorname{index}(D_{12}),\qquad
  \operatorname{index}(D_1)=\operatorname{index}(D_{02}).
\]
$D_0$ 是 $D_{01}$ 与 $D_{12}$ 的直和，由定理~\ref{thm:C-1-10}(i)
可知它是\BookRecordIndex{T0128}弗雷德霍姆算子。$D_1$ 与 $D_{02}$ 同构：定义域的同构将
$\xi_{02}\in W_{F_{02}}^{1,2}(\Sigma_{02},E_{02})$
送到 $(\xi_{02}|_{\Sigma_{01}},\xi_{02}|_{\Sigma_{12}})\in X_1$；
值域同理。（注意：这里取 $k=1$，因为 $k\ge2$ 时映射
$\xi_{02}\mapsto(\xi_{02}|_{\Sigma_{01}},\xi_{02}|_{\Sigma_{12}})$
对 $W^{k,2}$ 并非满射。）这证明了步骤 3。

\noindent\textbf{步骤 4：} 对 $t\in\mathbb R/\mathbb Z$ 定义全实子空间
$\Lambda(t)\subset\mathbb C^n$ 为
\[
  \Lambda(t):=\Phi(z_t)^{-1}F_{1,z_t},\qquad z_t:=\phi(0,t)\in U.
\]
则存在光滑映射
$\Psi:[0,1]\times\mathbb R/\mathbb Z\to
  \operatorname{End}_{\mathbb R}(\mathbb C^n\oplus\mathbb C^n)$，满足
\begin{equation}\label{eq:C-4-3}
  \Psi(s,t)I=I\Psi(s,t),\qquad
  I:=\begin{pmatrix}-i&0\\0&i\end{pmatrix},
\end{equation}
对所有 $s,t$ 成立，且
\begin{equation}\label{eq:C-4-4}
  \Psi(0,t)^{-1}\Delta=\Lambda(t)\oplus\Lambda(t),
\end{equation}
并在 $1/2\le s\le1$ 时 $\Psi(s,t)=\mathbf 1$。
在 $(\mathbb C^n\oplus\mathbb C^n,I)$ 中，全实子空间圈
$\Lambda_0(t):=\Lambda(t)\oplus\Lambda(t)$ 的马斯洛夫指标为零，
故与常值圈
\[
  \Lambda_1(t):=\Delta:=\{(\zeta,\zeta)\mid\zeta\in\mathbb C^n\}
\]
同伦。选光滑同伦
$[0,1]\times\mathbb R/\mathbb Z\to R(\mathbb C^n\oplus\mathbb C^n,I)$，
$(s,t)\mapsto\Lambda(s,t)$，使 $\Lambda(0,t)=\Lambda_0(t)$，
且 $1/2\le s\le1$ 时 $\Lambda(s,t)=\Lambda_1(t)$。
再选 $\Lambda(s,t)$ 的全局光滑\BookRecordIndex{T0067}标架
$e_1(s,t),\ldots,e_{2n}(s,t)$，使对所有 $i$ 及 $1/2\le s\le1$，
$e_i(s,t)=e_i(1,t)$。令
\[
  \Psi(s,t)e_i(s,t):=e_i(1,t),\qquad
  \Psi(s,t)Ie_i(s,t):=Ie_i(1,t)
\]
对所有 $i,s,t$ 成立。于是 $\Psi$ 满足式~\eqref{eq:C-4-3}、
\eqref{eq:C-4-4}，且 $1/2\le s\le1$ 时 $\Psi(s,t)=\mathbf1$。
步骤 4 得证。

\noindent\textbf{步骤 5：} 存在希尔伯特空间同构
$\Psi_X:X_0\to X_1$、$\Psi_Y:Y\to Y$，
使算子 $D_1\circ\Psi_X-\Psi_Y\circ D_0:X_0\to Y$ 紧。
设 $\Psi$ 如步骤 4，并写
\[
  \Psi(s,t)=:
  \begin{pmatrix}A(s,t)&B(s,t)\\C(s,t)&D(s,t)\end{pmatrix}.
\]
由式~\eqref{eq:C-4-3}，$A(s,t),D(s,t)$ 为复线性自同态，
在 $s=1$ 附近均为 $\mathbf1$；$B(s,t),C(s,t)$ 为复反线性，
在 $s=1$ 附近均为零。定义 $\Psi_X:X_0\to X_1$ 为
$\Psi_X(\xi_{01},\xi_{12}):=(\widetilde\xi_{01},\widetilde\xi_{12})$，
其中 $\widetilde\xi_{ij}=\xi_{ij}$ 于 $\Sigma_{ij}\setminus U$，
而在 $\Sigma_{ij}\cap U$ 上，由下式定义：
\begin{equation}\label{eq:C-4-5}
  \widetilde\xi_{ij}(\phi(s,t))
  :=\Phi(\phi(s,t))\widetilde\zeta_{ij}(s,t),\qquad
  \xi_{ij}(\phi(s,t))=:\Phi(\phi(s,t))\zeta_{ij}(s,t),
\end{equation}
其中
\begin{equation}\label{eq:C-4-6}
\begin{aligned}
  \widetilde\zeta_{01}(s,t)
    &:=A(-s,t)\zeta_{01}(s,t)+B(-s,t)\zeta_{12}(-s,t),
      &&-1\le s\le0,\\
  \widetilde\zeta_{12}(s,t)
    &:=C(s,t)\zeta_{01}(-s,t)+D(s,t)\zeta_{12}(s,t),
      &&0\le s\le1.
\end{aligned}
\end{equation}
若 $(\xi_{01},\xi_{12})\in X_0$，则每个 $z\in\Gamma_1$ 上
$\xi_{01}(z),\xi_{12}(z)\in F_{1,z}$，因而
$(\zeta_{01}(0,t),\zeta_{12}(0,t))\in\Lambda(t)\oplus\Lambda(t)$。
故 $(\widetilde\zeta_{01}(0,t),\widetilde\zeta_{12}(0,t))\in\Delta$，
从而 $\widetilde\xi_{01}(z)=\widetilde\xi_{12}(z)$，$z\in\Gamma_1$。
这表明 $\Psi_X$ 将 $X_0$ 送入 $X_1$。

再定义 $\Psi_Y:Y\to Y$ 为
$\Psi_Y(\eta_{01},\eta_{12}):=(\widetilde\eta_{01},\widetilde\eta_{12})$，
其中 $\widetilde\eta_{ij}=\eta_{ij}$ 于 $\Sigma_{ij}\setminus U$，
而在 $\Sigma_{ij}\cap U$ 上，由下式定义：
\begin{equation}\label{eq:C-4-7}
  \phi^*\widetilde\eta_{ij}:=(\Phi\circ\phi)\widetilde\beta_{ij},\qquad
  \phi^*\eta_{ij}=:(\Phi\circ\phi)\beta_{ij}(s,t),
\end{equation}
其中
\begin{equation}\label{eq:C-4-8}
\begin{aligned}
  \widetilde\beta_{01}(s,t)
    &:=A(-s,t)\beta_{01}(s,t)-B(-s,t)\beta_{12}(-s,t),
      &&-1\le s\le0,\\
  \widetilde\beta_{12}(s,t)
    &:=-C(s,t)\beta_{01}(-s,t)+D(s,t)\beta_{12}(s,t),
      &&0\le s\le1.
\end{aligned}
\end{equation}
取 $(\xi_{01},\xi_{12})\in X_0$，并按式~\eqref{eq:C-4-5}、
\eqref{eq:C-4-6} 定义 $\zeta_{ij}(s,t),\widetilde\zeta_{ij}(s,t)\in\mathbb C^n$。
由于 $B$ 复反线性，链式法则给出下列两式。
\sourcefootnote{原书实体页 617 在第一式的最后一项把
$\zeta_{12}(-s,t)$ 写成 $\zeta_{12}(s,t)$，与式 (C.4.6) 的反射自变量不符；
在第二式把系数导数 $(\partial_sC+i\partial_tC)$ 前误置负号。
两处按式 (C.4.6) 直接求导补正，未改变其余项。}
当 $-1\le s\le0$ 时，
\begin{align*}
  (\partial_s+i\partial_t)\widetilde\zeta_{01}(s,t)
  ={}&A(-s,t)(\partial_s+i\partial_t)\zeta_{01}(s,t)\\
    &-B(-s,t)(\partial_s+i\partial_t)\zeta_{12}(-s,t)\\
    &+(-\partial_sA+i\partial_tA)(-s,t)\zeta_{01}(s,t)\\
    &+(-\partial_sB+i\partial_tB)(-s,t)\zeta_{12}(-s,t);
\end{align*}
当 $0\le s\le1$ 时，
\begin{align*}
  (\partial_s+i\partial_t)\widetilde\zeta_{12}(s,t)
  ={}&-C(s,t)(\partial_s+i\partial_t)\zeta_{01}(-s,t)\\
    &+D(s,t)(\partial_s+i\partial_t)\zeta_{12}(s,t)\\
    &+(\partial_sC+i\partial_tC)(s,t)\zeta_{01}(-s,t)\\
    &+(\partial_sD+i\partial_tD)(s,t)\zeta_{12}(s,t).
\end{align*}
但由步骤 1、2，$D_1(\Psi_X\xi)=\Phi(\phi)\bar\partial\widetilde\zeta$。
故据 $\Psi_Y$ 的定义，$D_1\Psi_X-\Psi_YD_0$ 为紧算子。
步骤 5 得证。

由步骤 5 与定理~\ref{thm:A-1-5}，$D_0,D_1$
具有相同的\BookRecordIndex{T0129}弗雷德霍姆指标，再由步骤 3 即得结论。
这证明了定理~\ref{thm:C-4-2}。
\end{proof}

\begin{proof}[定理~\ref{thm:C-1-10} 的证明]
断言 (i) 已在 C.2 节证明。断言 (ii) 的指标公式基于定理~\ref{thm:C-4-2}。
由式~\eqref{eq:C-4-2} 及\BookRecordIndex{T0062}边界马斯洛夫指标的复合公理，
若指标公式对分解中的三个曲面 $\Sigma_{ij}$ 中的两个成立，
则它对第三个也成立。因此，断言 (ii) 由定理~\ref{thm:C-4-1}
及 C.3 节的裤形归纳\BookRecordIndex{T0326}推出。

下面证明断言 (iii)。第一步，将 $\mu(E,F)<0$ 时的单射性结论，
通过把 $\Sigma$ 与其反向曲面\BookRecordIndex{T0254}粘合，归结为闭黎曼曲面的情形。
具体如下：设 $\Sigma$ 是带非空边界 $\partial\Sigma=\Gamma$
的紧致连通黎曼曲面，带有平凡丛 $\Sigma\times\mathbb C$；
设 $F\subset\Gamma\times\mathbb C$ 是边界上的全实子丛。
则存在函数 $\lambda:\Gamma\to S^1/\{\pm1\}$，使
\[
  F_z=\lambda(z)\mathbb R,\qquad z\in\Gamma.
\]
截面 $\xi:\Sigma\to\mathbb C$ 满足 $\xi(z)\in F_z$，
当且仅当
\[
  \overline{\xi(z)}=\lambda(z)^{-2}\xi(z),\qquad z\in\Gamma.
\]
考虑闭黎曼曲面
\[
  S:=\Sigma\times\{0,1\}/\sim,
\]
其中 $\Sigma\times\{1\}$ 上的复结构反向，且对 $z\in\Gamma$，
$(z,0)\sim(z,1)$。令 $\gamma:=\lambda^{-2}:\Gamma\to S^1$，
并考虑线丛
\[
  E:=E_0\times_\gamma E_1\longrightarrow S,
\]
其中 $E_0:=(\Sigma\times\{0\})\times\mathbb C$、
$E_1:=(\Sigma\times\{1\})\times\mathbb C$ 是平凡丛，
且对 $z\in\Gamma$、$\zeta\in\mathbb C$，
$(z,0,\zeta)\sim(z,1,\gamma(z)\zeta)$。
$E$ 的一个截面 $\zeta:S\to E$ 由函数对
$\zeta_0,\zeta_1:\Sigma\to\mathbb C$ 给出，满足
\[
  \zeta_1(z)=\gamma(z)\zeta_0(z),\qquad z\in\Gamma.
\]
由定理~\ref{thm:C-3-10}，原书写出 $E$ 的陈数为
\sourcefootnote{原书把反向曲面上的第二个\BookRecordIndex{T0062}边界马斯洛夫指标直接写成与第一个相同。
按 C.3 节从曲面诱导的边界定向定义，若第二份纤维仍取原复结构和同一平凡化，
其指标应变号；这里未明说为抵消该变号所需的共轭纤维约定。
此处保留原式，关于\BookRecordIndex{T0254}粘合后复线丛及陈数的约定需在复核时另行判定。}
\[
  2\langle c_1(E),[S]\rangle
   =\mu(E_0,F)+\mu(E_1,F)
   =2\mu(\Sigma\times\mathbb C,F)<0.
\]
要在光滑范畴中使这一构造精确，还须选取跨越边界曲线 $C$
的 $S$ 的坐标卡及 $E$ 的局部平凡化；$E$ 的光滑截面于是还须
在 $C$ 上满足各阶高导数的附加\BookRecordIndex{T0254}粘合条件。
不过，只处理 $W^{1,p}$ 截面和 $L^p$ 类\BookRecordIndex{T0210}柯西--黎曼算子时，
无需作这些选择。

现考虑 $\Sigma$ 上形如 $\bar\partial+\alpha$ 的\BookRecordIndex{T0210}柯西--黎曼算子，
其中对某个 $p>2$，
\[
  \alpha\in L^p(\Sigma,\Lambda^{0,1}T^*\Sigma
    \otimes\operatorname{End}_{\mathbb R}(\mathbb C)).
\]
此算子确定 $E$ 上一个 $L^p$ 类\BookRecordIndex{T0210}柯西--黎曼算子 $D$：
在 $\Sigma\times\{0\}$ 上为 $\bar\partial+\alpha$，
在 $\Sigma\times\{1\}$ 上为 $\partial+\bar\alpha$。这里
\[
  \bar\alpha(z,\widehat z):=
    \tau\circ\alpha(z,\widehat z)\circ\tau,
  \qquad \widehat z\in T_z\Sigma,
\]
其中 $\tau:\mathbb C\to\mathbb C$ 为复共轭。
若函数 $\xi\in W^{1,p}(\Sigma,\mathbb C)$ 满足
$\bar\partial\xi+\alpha\xi=0$ 及 $\xi(\partial\Sigma)\subset F$，
则由 $\zeta_0(z):=\xi(z)$、$\zeta_1(z):=\overline{\xi(z)}$
给出 $D$ 核中的截面 $\zeta\in W^{1,p}(S,E)$。
故只须在闭曲面上证明 $L^p$ 类\BookRecordIndex{T0210}柯西--黎曼算子的断言 (iii)。

现在令 $E\to S$ 是闭黎曼曲面 $S$ 上的复线丛，$D$ 是 $S$ 上的
$L^p$ 类\BookRecordIndex{T0210}柯西--黎曼算子，其中 $p>2$。我们要证明当 $\mu(E)<0$ 时
$D$ 单射。先考虑 $D$ 复线性的情形。由习题~\ref{ex:C-1-8}，
$D$ 经规范变换等价于光滑的复线性\BookRecordIndex{T0210}柯西--黎曼算子 $D_0$。
由注~\ref{rem:C-1-1}，$E$ 上存在使 $D_0$ 成为通常的 $\bar\partial$
算子的\BookRecordIndex{T0288}全纯结构。因此，$D_0$ 核中的任意截面 $\xi:S\to E$
在局部都可表示为开集上的全纯函数。若 $\xi$ 非零，则其零点孤立，
且指标为正。由于 $E$ 的陈数是仅有孤立零点的截面的各零点指标之和，
当 $\mu(E)=2c_1(E)<0$ 时，$D_0$ 的核只能为零。
$D$ 与 $D_0$ 的核同构，故定理~\ref{thm:C-1-10}(iii) 的第一个断言
对复线性\BookRecordIndex{T0210}柯西--黎曼算子成立。

处理实线性情形时，我们采用 Hofer--Lizan--Sikorav~\srcite{179}
的一个技巧。（定理~\ref{thm:2-3-5} 证明的第二步也用了同样的思路。）
在复线丛 $E\to S$ 上选择光滑的复线性\BookRecordIndex{T0210}柯西--黎曼算子 $D_0$，并写为
\[
  D=D_0+a,\qquad
  a\in L^p\bigl(S,\Lambda^{0,1}T^*S
       \otimes\operatorname{End}_{\mathbb R}(E)\bigr).
\]
令 $\xi\in W^{1,p}(S,E)$\sourcefootnote{原书在闭曲面 $S$ 上仍写
$W_F^{1,p}(S,E)$，但此时并无边界条件 $F$；据定义域改为
$W^{1,p}(S,E)$。} 满足 $D\xi=D_0\xi+a\xi=0$。
由于 $\xi$ 连续，而 $E$ 中两个元素的比是良定义的复数，
可以如下定义 $b\in L^p(S,\Lambda^{0,1}T^*S)$：
\[
  b(z,\widehat z):=
  \begin{cases}
    \bigl(a(z,\widehat z)\xi(z)\bigr)/\xi(z),&\xi(z)\ne0,\\
    0,&\text{其他情形},
  \end{cases}
  \qquad \widehat z\in T_zS.
\]
于是 $D_0+b$ 是复线性\BookRecordIndex{T0210}柯西--黎曼算子，且由构造知
$a\xi=b\xi$，从而 $D_0\xi+b\xi=0$。换言之，实线性\BookRecordIndex{T0210}柯西--黎曼
算子核中的每个元素，也属于同一线丛上的某个 $L^p$ 类复线性
\BookRecordIndex{T0210}柯西--黎曼算子的核。由刚证得的结论，负陈数复线丛上的
每个实线性\BookRecordIndex{T0210}柯西--黎曼算子均为单射。

这就证明了定理~\ref{thm:C-1-10}(iii) 的第一个断言在一般情形下成立。
为证明第二个断言，注意由定理~\ref{thm:C-1-10}(i) 的推论~\ref{cor:C-1-11}，
$D$ 的余核同构于\BookRecordIndex{T0080}丛对
\[
  \bigl(\Lambda^{1,0}T^*\Sigma\otimes_{\mathbb C}E^*,
       T\partial\Sigma\otimes_{\mathbb R}(E/F)^*\bigr)
\]
上某个\BookRecordIndex{T0210}柯西--黎曼算子的核；该\BookRecordIndex{T0080}丛对的\BookRecordIndex{T0062}边界马斯洛夫指标为
\[
  \mu\bigl(\Lambda^{1,0}T^*\Sigma\otimes_{\mathbb C}E^*,
       T\partial\Sigma\otimes_{\mathbb R}(E/F)^*\bigr)
  =-\mu(E,F)-2\chi(\Sigma).
\]
故当 $\mu(E,F)+2\chi(\Sigma)>0$ 时，余核消失。
这证明了定理~\ref{thm:C-1-10}。
\end{proof}
